\documentclass{article}
\usepackage{graphicx} % Required for inserting images
\usepackage{booktabs, tabularx}
\usepackage{pdflscape}
\usepackage{xltabular}
\usepackage{pdflscape}
\usepackage[margin=2.5cm]{geometry}

\title{Interface Specification}
\author{Frédéric Kah}
\date{October 2026}

\begin{document}

\maketitle

\section{Introduction}

\begin{landscape}
\begin{table}[p]
  \centering
  \small
  \renewcommand{\arraystretch}{1.3}
  \begin{tabularx}{\linewidth}{@{} l l >{\raggedright\arraybackslash}X >{\raggedright\arraybackslash}p{5cm} @{}}
    \toprule
    \textbf{Block} & \textbf{Owner} & \textbf{What it does} & \textbf{Talks to} \\
    \midrule
    \texttt{BoardTop} & Frédéric
      & Maps the design to the FPGA pins, synchronizes and inverts the reset button, drives the Debug LEDs
      & FPGA pins $\leftrightarrow$ \texttt{Tpu} \\
    \texttt{Tpu} (core top) & Frédéric
      & Instantiates and wires all blocks together; target of the full-system tests
      & All blocks \\
    \texttt{UartRx}, \texttt{UartTx} & Frédéric
      & Translate between the USB-serial link, where bytes travel one bit at a time,
    and the rest of the design, which works with whole bytes.
      & Internal to Frédéric's part \\
    \texttt{RxFifo} & Frédéric
      & Waiting line for received bytes: stores each byte from \texttt{UartRx}
    until the decoder is ready to read it, since the UART receiver cannot wait.
      & \texttt{UartRx}, decoder \\
    \texttt{CommandDecoder} & Frédéric
      & Parses the byte stream from
    \texttt{UartRx} into commands. Data transfers it handles itself: matrices
    sent by the host are written into the buffer (\texttt{WRITE}), and results
    are read from it (\texttt{READ}). Compute instructions (\texttt{MATMUL})
    are passed on to the sequencer, which executes them. As the only block
    connected to \texttt{UartTx}, it also sends back to the host everything the
    host asks for: results from the buffer and cycle counts from the
    sequencer's performance counters. Possible commands are \texttt{WRITE}, \texttt{READ}, \texttt{MATMUL}, \texttt{STATUS}.
      & UART, buffer, sequencer \\
    \texttt{UnifiedBuffer} & Giovanni
      & Banked on-chip memory holding the matrices A, B and the results C
      & Decoder, data mover, accumulators \\
    \texttt{Sequencer} (+ tiling) & Giovanni
      & Executes the compute instructions received from the
    decoder. A matrix product is usually larger than the array, so it splits it
    into array-sized pieces (tiles) and processes them one after another: for
    each tile, it tells the data mover which weights and rows to send into the
    array, and tells the accumulators whether to start a new sum or add to the
    previous one. It signals when it is busy, and counts how many cycles each
    instruction takes, so the decoder can report the performance to the host.
      & Decoder, data mover, accumulators \\
    \texttt{DataMover} & Giovanni
      & Reads weight and activation rows from the buffer and streams them into the array
      & Buffer, \textbf{array} \\
    \texttt{SystolicArray} (+ skew/de-skew) & Andreea
      & Grid of PEs computing one tile; skews inputs and de-skews the output columns so each output row appears in one cycle
      & Data mover, accumulators \\
    \texttt{Accumulators} & Giovanni
      & Add partial sums across K-tiles. On the last K-tile, applies the post-processing (optional requantization and ReLU) and writes the final results to the result memory
      & \textbf{Array}, buffer \\
    \bottomrule
  \end{tabularx}
  \caption{High-level blocks of the design, their owners and connections.}
  \label{tab:blocks}
\end{table}
\end{landscape}

% ============================================================
% Port tables for each block
% Preamble needs: \usepackage{booktabs, tabularx}
% ============================================================

\newcolumntype{Y}{>{\raggedright\arraybackslash}X}
% Common column layout: Port | Dir | Type | Proto | Connected to | Description
\newcommand{\porthead}{%
  \toprule
  \textbf{Port} & \textbf{Dir} & \textbf{Type} & \textbf{Proto} & \textbf{Connected to} & \textbf{Description} \\
  \midrule}

\subsection*{Conventions}

\paragraph{Reading the tables.}
\textbf{Dir} is given from the block's own point of view (\emph{in} = the block receives it).
\texttt{inT} and \texttt{accT} are the input and accumulator types defined in \texttt{SAConfig}
(int8 and int32 in the MVP); \texttt{rows} and \texttt{cols} are the array dimensions.
Ports marked $\dagger$ are internal to Giovanni's part; their exact shape is his choice.

\paragraph{Clock.}
The whole design uses a single clock, the 100\,MHz board clock. There is no second clock domain,
so no signal ever needs to be synchronized between blocks. The only asynchronous inputs are the
UART line and the reset button, which are synchronized where they enter the design.

\paragraph{Reset.}
Reset is \textbf{synchronous and active-high}: it takes effect on a rising clock edge while
\texttt{reset} is 1, and one cycle is enough. Only \texttt{BoardTop} sees the board's active-low
button; it synchronizes and inverts it, and every other block receives the same active-high reset.
The general rule, refined in each table, is:
\begin{itemize}
  \item every block returns to its idle state, and any half-finished command, transfer or tile is discarded;
  \item every \texttt{valid} output is 0, so no transfer can happen in the cycle after reset;
  \item \textbf{memory contents are not cleared}: block RAM cannot be reset in one cycle, so after a reset
        the host must assume the buffer holds garbage and write its matrices again;
  \item data registers that are always written before being read (weights, partial sums) need no reset,
        which saves logic; only control state and \texttt{valid} bits are reset.
\end{itemize}

\subsection*{Protocols}

Two blocks exchange information in one of three ways, written in the \textbf{Proto} column.

\subsubsection*{RV: ready/valid (Chisel \texttt{Decoupled})}

Used when the receiver may be busy or slower than the sender, so it must be able to say ``not now''.

\begin{itemize}
  \item \textbf{Signals.} \texttt{bits} (the data) and \texttt{valid} go from sender to receiver;
        \texttt{ready} goes from receiver to sender.
        \texttt{valid} = 1 means ``\texttt{bits} holds an item right now''; \texttt{ready} = 1 means ``I can take an item this cycle''.
  \item \textbf{Transfer.} One item is transferred on every rising clock edge where \texttt{valid} and \texttt{ready}
        are both 1 (Chisel calls this \texttt{fire}). In every other cycle, nothing is transferred.
  \item \textbf{Sender rules.} Once \texttt{valid} is 1, it stays 1 and \texttt{bits} stays unchanged until the
        transfer happens: an offered item cannot be withdrawn. \texttt{valid} must never depend on \texttt{ready}
        (otherwise each side waits for the other and the design deadlocks or forms a combinational loop).
  \item \textbf{Receiver rules.} \texttt{ready} may be set at any time, even before \texttt{valid}, and may depend on \texttt{valid}.
        Holding \texttt{ready} at 0 makes the sender wait (\emph{backpressure}).
  \item \textbf{In Chisel.} \texttt{Decoupled(T)} on the sender, \texttt{Flipped(Decoupled(T))} on the receiver.
\end{itemize}

\begin{center}\small
\begin{tabular}{l ccccc}
  \toprule
  Cycle            & 0 & 1 & 2 & 3 & 4 \\
  \midrule
  \texttt{valid}   & 0 & 1 & 1 & 1 & 0 \\
  \texttt{ready}   & 1 & 0 & 0 & 1 & 1 \\
  \texttt{bits}    & -- & X & X & X & -- \\
  transfer         & -- & -- & -- & \textbf{X} & -- \\
  \bottomrule
\end{tabular}
\end{center}
The sender offers X in cycle 1. The receiver is busy in cycles 1 and 2, so the sender keeps X on \texttt{bits};
X is transferred in cycle 3, the first cycle where both signals are 1.

\subsubsection*{V: valid-only (Chisel \texttt{Valid})}

Used when the sender cannot wait: the UART receiver (bits arrive on the wire whenever the PC sends them)
and the systolic array (every value moves one step per cycle, at times fixed by the schedule).

\begin{itemize}
  \item \textbf{Signals.} \texttt{bits} and \texttt{valid}, both from sender to receiver. There is no \texttt{ready}.
  \item \textbf{Transfer.} Every cycle where \texttt{valid} = 1 transfers one item. The receiver must take it in that
        same cycle; there is no second chance, and an item it does not take is lost.
  \item \textbf{Sender rules.} \texttt{valid} is 1 for exactly one cycle per item. Keeping it at 1 for two cycles
        means sending two items.
  \item \textbf{Receiver rules.} It must always be able to accept. If it may occasionally be busy, a FIFO goes in
        front of it (as behind \texttt{UartRx}).
  \item \textbf{Gaps.} A cycle with \texttt{valid} = 0 is simply ``no item''. Through the array, the \texttt{valid} bit
        travels alongside the data, so a gap at the input comes out as a gap at the output.
  \item \textbf{In Chisel.} \texttt{Valid(T)} on the sender, \texttt{Flipped(Valid(T))} on the receiver.
\end{itemize}

\begin{center}\small
\begin{tabular}{l ccccc}
  \toprule
  Cycle            & 0 & 1 & 2 & 3 & 4 \\
  \midrule
  \texttt{valid}   & 0 & 1 & 1 & 0 & 1 \\
  \texttt{bits}    & -- & A & B & -- & C \\
  transfer         & -- & \textbf{A} & \textbf{B} & -- & \textbf{C} \\
  \bottomrule
\end{tabular}
\end{center}

\subsubsection*{W: plain wire}

Used for information that is always meaningful, which the reader looks at whenever it needs it.

\begin{itemize}
  \item \textbf{Signals.} Just the value, with no \texttt{valid} or \texttt{ready}. There is no notion of transfer.
  \item \textbf{Level.} A value held for as long as it is true, such as \texttt{busy}, or a counter that the decoder reads
        when the host sends \texttt{STATUS}.
  \item \textbf{Pulse.} A Bool that is 1 for exactly one cycle to signal an event, such as \texttt{done}. A pulse is
        always called one in its description, so it is never mistaken for a level.
\end{itemize}

\paragraph{Performance counters.}
\texttt{cycles}: clock cycles from the acceptance of a \texttt{MATMUL} until \texttt{busy} drops (UART transfer time excluded).
\texttt{busyPeCycles}: PE-cycles in which a PE multiplies valid data, computed by the sequencer as $m \cdot k \cdot n$ (padded sizes), so no extra port on the array is needed. Utilization $= \texttt{busyPeCycles} / (\texttt{cycles} \cdot \texttt{rows} \cdot \texttt{cols})$.
Both are cleared when an instruction is accepted and hold their value after it finishes; reset also clears them.

\paragraph{Derived values and assumptions (MVP).}
\begin{itemize}
  \item $V = 1$: \texttt{wIn} and \texttt{aIn} are written for one product per PE per cycle. For $V > 1$ both vectors become $V$ times longer and a tile covers $\texttt{rows}\cdot V$ values of K.
  \item \texttt{rows} = \texttt{cols} = \texttt{banks}, a power of two (in practice 4 or 8: a 16$\times$16 array would need 256 DSPs, more than the 240 available).
  \item Matrices are row-major. Element $(i,j)$ of a matrix with base address $b$ and leading dimension
        \texttt{ld} sits at $b + i\cdot\texttt{ld} + j$. The host sends \texttt{ld} in \texttt{MATMUL}:
        \texttt{ldA} and \texttt{ldB} in the input memory, \texttt{ldC} in the result memory. For a dense
        matrix, \texttt{ld} is its number of columns; a larger \texttt{ld} selects a block of a bigger matrix.
        The host guarantees $\texttt{ldA} \ge k$, $\texttt{ldB} \ge n$, $\texttt{ldC} \ge n$;
        \texttt{aAddr}, \texttt{bAddr}, \texttt{ldA}, \texttt{ldB} are multiples of \texttt{banks};
        \texttt{cAddr} and \texttt{ldC} are multiples of \texttt{cols}. Violations are undefined behaviour.
\end{itemize}
\subsection*{\texttt{SAConfig} (Frédéric)}

\texttt{SAConfig} is a Scala case class holding every generator parameter. Each module receives it as a
constructor argument (\texttt{new SystolicArray(c)}) and takes all its sizes and types from it: nobody
hardcodes a width or a size anywhere else. The team chooses the \emph{base parameters}
(Table~\ref{tab:base}); the \emph{derived values} (Table~\ref{tab:derived}) are computed from them and
must never be set by hand. The column \textbf{In the MVP} says how far each one is actually exercised:
whether it can really vary, or whether only one value is accepted for now.

\begin{table}[htbp]
\small\renewcommand{\arraystretch}{1.25}
\begin{tabularx}{\linewidth}{@{} l l l Y Y @{}}
  \toprule
  \textbf{Parameter} & \textbf{Type} & \textbf{Default} & \textbf{Meaning} & \textbf{In the MVP} \\
  \midrule
  \texttt{rows}     & Int     & 4    & Number of PE rows. One tile covers \texttt{rows} values of K.
                    & \textbf{Varies}: 1 (single-MAC baseline), 2, 4, 8 for the scaling plots. Must equal \texttt{cols}. \\
  \texttt{cols}     & Int     & 4    & Number of PE columns. One tile produces \texttt{cols} columns of C.
                    & \textbf{Varies} with \texttt{rows} (square arrays only). \\
  \texttt{vecWidth} & Int     & 1    & Products each PE computes per cycle ($V$). A tile then covers $\texttt{rows}\cdot V$ values of K.
                    & \textbf{Fixed to 1}; other values are rejected. The ports are written for $V = 1$. \\
  \texttt{inWidth}  & Int     & 8    & Bits of an input element of A and B.
                    & \textbf{Fixed to 8}: the host protocol sends one byte per element. \\
  \texttt{signed}   & Boolean & true & Whether inputs and accumulators are signed.
                    & \textbf{Only true is tested}: the quantization scheme is symmetric signed int8. \\
  \texttt{accWidth} & Int     & 32   & Bits of an accumulator and of a raw result.
                    & \textbf{Fixed to 32}: \texttt{READ} sends at most 4 bytes per element. \\
  \texttt{dataflow} & enum    & \texttt{WeightStationary} & What stays inside the PEs while the other data flows through.
                    & \textbf{Only \texttt{WeightStationary}}; the parameter exists so other dataflows can be added later. \\
  \texttt{loadMode} & enum    & \texttt{Space} & How weights are loaded: \texttt{Space} shifts one row in per cycle (few wires), \texttt{Speed} loads all rows at once (many wires).
                    & \textbf{Only \texttt{Space}} (shift-in loading). \\
  \texttt{inDepth}  & Int     & 4096 & Number of \texttt{inT} elements the input memory holds (A and B together).
                    & \textbf{Varies}: any multiple of \texttt{banks} up to $2^{16}$, limited by the FPGA's block RAM. \\
 \texttt{outDepth} & Int & 4096 & Number of \texttt{accT} words the result memory holds (C). Separate from the
                                   input memory: results are 32 bits wide and are read back by the host one by one.
                    & \textbf{Varies}: any multiple of \texttt{cols} up to $2^{16}$. \\
  \texttt{maxM}     & Int     & 256  & Rows of C the accumulators can hold, i.e.\ the largest \texttt{m} one \texttt{MATMUL} accepts. Larger products are split into several \texttt{MATMUL}s by the host.
                    & \textbf{Varies}. The host reads it to check $m \le \texttt{maxM}$. \\
  \bottomrule
\end{tabularx}
\caption{Base parameters of \texttt{SAConfig}.}
\label{tab:base}
\end{table}

\begin{table}[htbp]
\small\renewcommand{\arraystretch}{1.25}
\begin{tabularx}{\linewidth}{@{} l >{\raggedright\arraybackslash}p{3.3cm} Y Y @{}}
  \toprule
  \textbf{Value} & \textbf{Definition} & \textbf{Meaning and users} & \textbf{In the MVP} \\
  \midrule
  \texttt{inT}  & \texttt{SInt(inWidth.W)} (\texttt{UInt} if unsigned)
                & Chisel type of an input element. Used by the decoder, buffer, data mover and array.
                & Always \texttt{SInt(8.W)}. \\
  \texttt{accT} & \texttt{SInt(accWidth.W)} (\texttt{UInt} if unsigned)
                & Chisel type of a partial sum or result. Used by the array, accumulators, buffer and decoder.
                & Always \texttt{SInt(32.W)}. \\
  \texttt{banks} & \texttt{rows} ($=$ \texttt{cols})
                & Number of input-memory banks, i.e.\ elements the data mover reads per cycle. Element address $a$ is in bank $a \bmod \texttt{banks}$, row $a / \texttt{banks}$.
                & Follows \texttt{rows}. \\
\texttt{bankBits} & $\log_2 \texttt{banks}$
                & Bits of an element address that select the bank: the low \texttt{bankBits} bits give the bank,
                  the others the row inside it. Used by the buffer to split addresses.
                & 2 with the defaults; 0 for the $1\times1$ baseline (one bank, no bank bits). \\
  \texttt{addrW} & $\lceil\log_2 \texttt{inDepth}\rceil$
                & Bits of an element address in the input memory (\texttt{aAddr}, \texttt{bAddr}, \texttt{WRITE} addresses).
                & 12 with the defaults. \\
  \texttt{rowAddrW} & $\texttt{addrW} - \log_2 \texttt{banks}$
                & Bits of a row address inside a bank. Used by the data mover's row reads (\texttt{BufRowReq}).
                & 10 with the defaults. \\
\texttt{outAddrW} & $\lceil\log_2 \texttt{outDepth}\rceil$
                & Bits of an element address in the result memory (\texttt{cAddr}, \texttt{READ} addresses).
                & 12 \\
  \texttt{shiftW} & $\lceil\log_2 \texttt{accWidth}\rceil$
                & Bits of the requantization shift: enough to shift by up to $\texttt{accWidth}-1$.
                  Used by \texttt{Instruction} and the requantization.
                & 5 \\
  \texttt{dimW} & 16 (constant)
                & Bits of the fields \texttt{m}, \texttt{k}, \texttt{n} and the leading dimensions \texttt{ld}. Fixed by the 2-byte fields of the host protocol, not computed.
                & Always 16. \\
  \texttt{arrayLatency} & $\texttt{rows} + \texttt{cols} - 1$ (to confirm with Andreea)
                & Cycles between row $m$ of \texttt{aIn} entering the array and row $m$ of \texttt{out} leaving it, after skew and de-skew. Comes from Andreea's schedule function; used by the sequencer and accumulators.
                & 7 for a $4\times4$ array, if each PE hop adds one register. \\
    \texttt{weightToActGap} & 0 (to confirm with Andreea)
                & Minimum number of empty cycles between the last weight row on \texttt{wIn} and the first
                  activation row on \texttt{aIn}: the first activation may arrive $\texttt{weightToActGap}+1$
                  cycles after the last weight. Depends on the PE design; used by the data mover.
                & 0 \\
  \texttt{dspCount} & $\texttt{rows}\cdot\texttt{cols}\cdot\texttt{vecWidth}$
                & Number of multipliers, one DSP each. Used by the checks and the resource reports.
                & 16 for $4\times4$; 64 for $8\times8$. \\
  \bottomrule
\end{tabularx}
\caption{Derived values of \texttt{SAConfig}, computed from the base parameters.}
\label{tab:derived}
\end{table}

\paragraph{Checks.}
The constructor rejects any configuration the MVP does not support, with a \texttt{require} and a clear message:
\begin{itemize}
  \item \texttt{dataflow} is \texttt{WeightStationary}, \texttt{loadMode} is \texttt{Space}, \texttt{vecWidth} is 1;
  \item $\texttt{rows} = \texttt{cols}$, a power of two;
  \item $\texttt{dspCount} \le 240$ (DSP slices of the XC7A100T, so at most $8\times8$ in practice);
  \item \texttt{inDepth} is a multiple of \texttt{banks} and $\texttt{addrW} \le 16$ (2-byte addresses);
  \item $\texttt{accWidth} \ge 2\cdot\texttt{inWidth} + \texttt{dimW}$, so that a sum over any K the protocol can express never overflows.
\end{itemize}

\paragraph{Not in \texttt{SAConfig}.}
Board settings (\texttt{clockHz}, \texttt{baud}) and the \texttt{RxFifo} depth live in a separate \texttt{BoardConfig},
because they do not change the accelerator itself.


\subsection*{\texttt{BoardConfig} (Frédéric)}

\texttt{BoardConfig} holds the settings of the link between the PC and the FPGA. They are kept separate
from \texttt{SAConfig} because they do not change the accelerator itself: the same array, buffer and
sequencer work with any clock or baud rate. Only \texttt{UartRx}, \texttt{UartTx}, \texttt{RxFifo} and
\texttt{Tpu} receive it. Tests use a different \texttt{BoardConfig} from the board (a much higher baud
rate), while keeping the same \texttt{SAConfig}.

\begin{table}[htbp]
\small\renewcommand{\arraystretch}{1.25}
\begin{tabularx}{\linewidth}{@{} l l l Y Y @{}}
  \toprule
  \textbf{Parameter} & \textbf{Type} & \textbf{Default} & \textbf{Meaning} & \textbf{In the MVP} \\
  \midrule
  \texttt{clockHz}     & Int & 100\,000\,000 & Clock frequency in Hz. With \texttt{baud}, sets the length of one UART bit: $\texttt{clockHz}/\texttt{baud}$ cycles.
                       & \textbf{Fixed} to the Nexys4 DDR's 100\,MHz clock on the board. \\
  \texttt{baud}        & Int & 115\,200 & UART speed in bits per second. Must match the setting of the Python host.
                       & \textbf{Fixed} on the board; \textbf{much higher} in simulation (e.g.\ 4 cycles per bit) so full-system tests run fast. \\
  \texttt{rxFifoDepth} & Int & 16 & Number of bytes \texttt{RxFifo} can store while the decoder is busy. If it is full, the next byte is lost and \texttt{overrun} is set.
                       & \textbf{Fixed to 16}: enough for any normal stall of the decoder. \\
  \texttt{frameTimeout} & Int & $2^{20}$ & Cycles allowed between two bytes of a frame before the decoder drops the partial frame and sets \texttt{error}.
                       & \textbf{Fixed}; about 10\,ms at 100\,MHz. Smaller in simulation. \\
  \bottomrule
\end{tabularx}
\caption{Parameters of \texttt{BoardConfig}.}
\label{tab:board}
\end{table}

\paragraph{Checks.}
$\texttt{clockHz}/\texttt{baud} \ge 4$: the receiver samples each bit in its middle, so it needs several
cycles per bit. This also catches a mistyped baud rate in a test.

% ------------------------------------------------------------
\subsection*{\texttt{BoardTop} (Frédéric)}
{\small\renewcommand{\arraystretch}{1.25}
\noindent\begin{tabularx}{\linewidth}{@{} l l l l l Y @{}}
  \porthead
  \texttt{CLK100MHZ}     & in  & Clock     & W & FPGA pin        & 100\,MHz board clock. \\
  \texttt{CPU\_RESETN}   & in  & Bool      & W & FPGA pin        & Reset button, active-low and asynchronous. Synchronized with two flip-flops and inverted, then drives the \texttt{reset} of \texttt{Tpu}. \texttt{BoardTop} is a \texttt{RawModule}: it has no implicit clock or reset of its own. \\
  \texttt{UART\_TXD\_IN} & in  & Bool      & W & USB-UART bridge & Serial data from the PC (named from the PC's side, as in the Nexys \texttt{.xdc}). \\
  \texttt{UART\_RXD\_OUT}& out & Bool      & W & USB-UART bridge & Serial data to the PC. \\
  \texttt{LED}           & out & UInt(16)  & W & FPGA pins       & Debug indicators: busy, overrun, error. \\
  \bottomrule
\end{tabularx}}

% ------------------------------------------------------------
\subsection*{\texttt{Tpu} (core top, Frédéric)}
{\small\renewcommand{\arraystretch}{1.25}
\noindent\begin{tabularx}{\linewidth}{@{} l l l l l Y @{}}
  \porthead
  \texttt{clock} & in & Clock & -- & \texttt{BoardTop} & 100\,MHz clock shared by all blocks. \\
  \texttt{reset} & in & Bool & W & \texttt{BoardTop} & Synchronous, active-high. Returns the whole design to idle (see each block below). The host must resend any command that was in progress, and its matrices. \\
  \midrule
  \texttt{rxd}   & in  & Bool               & W & \texttt{BoardTop} & Serial data from the host. \\
  \texttt{txd}   & out & Bool               & W & \texttt{BoardTop} & Serial data to the host. \\
  \texttt{debug} & out & \texttt{DebugInfo} & W & \texttt{BoardTop} & Busy, overrun and error, shown on the LEDs.
     Unlike the \texttt{STATUS} flags, the LED copies of overrun and error stay on until reset. \\
  \bottomrule
\end{tabularx}}

% ------------------------------------------------------------
\subsection*{\texttt{UartRx} (Frédéric)}
Parameters: \texttt{clockHz}, \texttt{baud}.

{\small\renewcommand{\arraystretch}{1.25}
\noindent\begin{tabularx}{\linewidth}{@{} l l l l l Y @{}}
  \porthead
  \texttt{clock} & in & Clock & -- & \texttt{Tpu} & 100\,MHz clock shared by all blocks. \\
  \texttt{reset} & in & Bool & W & \texttt{Tpu} & Synchronous, active-high. Back to waiting for a start bit; a partly received byte is discarded; \texttt{out.valid} = 0. \\
  \midrule
  \texttt{rxd}        & in  & Bool    & W & \texttt{Tpu.rxd} & Serial line, idle high. Synchronized internally (two flip-flops). \\
  \texttt{out}        & out & UInt(8) & V & \texttt{RxFifo}          & Received byte, valid for one cycle. Cannot wait, hence the FIFO behind it. \\
  \texttt{framingErr} & out & Bool    & W & debug            & The stop bit was not 1: the byte is corrupted (usually a baud-rate mismatch). \\
  \bottomrule
\end{tabularx}}

% ------------------------------------------------------------
\subsection*{\texttt{UartTx} (Frédéric)}
Parameters: \texttt{clockHz}, \texttt{baud}.

{\small\renewcommand{\arraystretch}{1.25}
\noindent\begin{tabularx}{\linewidth}{@{} l l l l l Y @{}}
  \porthead
  \texttt{clock} & in & Clock & -- & \texttt{Tpu} & 100\,MHz clock shared by all blocks. \\
  \texttt{reset} & in & Bool & W & \texttt{Tpu} & Synchronous, active-high. Back to idle: \texttt{txd} = 1 (line idle), \texttt{in.ready} = 1. A byte being sent is cut short, so the PC may receive one corrupted byte. \\
  \midrule
  \texttt{in}  & in  & UInt(8) & RV & \texttt{CommandDecoder} & Byte to send. \texttt{ready} is high only while idle; one byte takes about 8\,700 cycles at 115\,200 baud. \\
  \texttt{txd} & out & Bool    & W  & \texttt{Tpu.txd}        & Serial line, idle high. \\
  \bottomrule
\end{tabularx}}

% ------------------------------------------------------------
\subsection*{\texttt{RxFifo} (Frédéric)}
A FIFO (\emph{first in, first out}) is a small memory that works like a waiting line: items
enter at the back and leave from the front, in arrival order. \texttt{RxFifo} sits between
\texttt{UartRx} and \texttt{CommandDecoder}. \texttt{UartRx} cannot wait (protocol V), while the
decoder sometimes must, for example when the buffer is busy or while it sends a reply through
\texttt{UartTx}. The FIFO stores the bytes that arrive in the meantime, so none are lost.
It is a standard Chisel \texttt{Queue(UInt(8.W), entries)}, instantiated in \texttt{Tpu}.
Parameter: \texttt{entries} (e.g.\ 16). No FIFO is needed on the transmit side, because the
decoder itself can wait for \texttt{UartTx}.

{\small\renewcommand{\arraystretch}{1.25}
\noindent\begin{tabularx}{\linewidth}{@{} l l l l l Y @{}}
  \porthead
  \texttt{clock} & in  & Clock   & -- & \texttt{Tpu}               & 100\,MHz clock shared by all blocks. \\
  \texttt{reset} & in  & Bool    & W  & \texttt{Tpu}               & Synchronous, active-high. The FIFO is emptied: bytes not yet read by the decoder are lost. \\
  \midrule
  \texttt{enq}   & in  & UInt(8) & RV & \texttt{UartRx.out}        & Bytes to store. \texttt{UartRx} ignores \texttt{ready}, so a byte that arrives while the FIFO is full
  (\texttt{enq.valid} = 1 and \texttt{enq.ready} = 0) is lost; this event, $\texttt{enq.valid} \land \lnot\texttt{enq.ready}$, is sent as a pulse to
\texttt{CommandDecoder.rxOverrun}. \\
  \texttt{deq}   & out & UInt(8) & RV & \texttt{CommandDecoder.rx} & Oldest stored byte; \texttt{valid} = 1 whenever the FIFO is not empty. \\
  \texttt{count} & out & UInt    & W  & (unused)                   & Number of stored bytes; available for debugging. \\
  \bottomrule
\end{tabularx}}
% ------------------------------------------------------------
\subsection*{\texttt{CommandDecoder} (Frédéric)}
{\small\renewcommand{\arraystretch}{1.25}
\noindent\begin{tabularx}{\linewidth}{@{} l l l l l Y @{}}
  \porthead
  \texttt{clock} & in & Clock & -- & \texttt{Tpu} & 100\,MHz clock shared by all blocks. \\
  \texttt{reset} & in & Bool & W & \texttt{Tpu} & Synchronous, active-high. Back to waiting for an opcode; a partly received command is discarded; all \texttt{valid} outputs = 0; \texttt{error} cleared. \\
  \midrule
  \texttt{rx}        & in  & UInt(8)                 & RV & \texttt{RxFifo}               & Bytes received from the host. \\
  \texttt{tx}        & out & UInt(8)                 & RV & \texttt{UartTx}        & Bytes to the host: status replies, results, counters. \\
  \texttt{bufWr}     & out & \texttt{BufWritePort}   & RV & \texttt{UnifiedBuffer} & Writes one matrix element (\texttt{WRITE} command). \\
  \texttt{bufRdReq}  & out & \texttt{BufReadReq}     & RV & \texttt{UnifiedBuffer} & Requests one result element (\texttt{READ} command). \\
  \texttt{bufRdResp} & in  & \texttt{accT}           & RV & \texttt{UnifiedBuffer} & The requested result. \\
  \texttt{instr}     & out & \texttt{Instruction}    & RV & \texttt{Sequencer}     & Compute instruction (\texttt{MATMUL}). Its acceptance starts the operation. \\
  \texttt{status}    & in  & \texttt{SeqStatus}      & W  & \texttt{Sequencer}     & Busy flag and performance counters, sent to the host on \texttt{STATUS}. \\
  \texttt{rxOverrun} & in & Bool & W & \texttt{RxFifo} & Pulse: a received byte was lost because \texttt{RxFifo} was full
 ($\texttt{enq.valid} \land \lnot\texttt{enq.ready}$). Sets the decoder's sticky \texttt{overrun} flag. \\
  \texttt{error}     & out & Bool                    & W  & debug                  & Unknown opcode or malformed command received. \\
  \bottomrule
\end{tabularx}}

\paragraph{Padding (MVP).}
The host pads A and B with zeros so that \texttt{k} and \texttt{n} are multiples of \texttt{banks}. Zeros do not change the sums. The instruction carries the padded sizes; the host discards the padded columns of C. \texttt{m} needs no padding but must not exceed \texttt{maxM} (accumulator depth). Violations are undefined behaviour: the Python host checks them before sending. (Later extension: zero masking in the data mover.)

\paragraph{Host protocol.}
Multi-byte fields are little-endian. Addresses, counts and dimensions are 2 bytes; one element is one byte (\texttt{inT} = int8).

\noindent\begin{tabularx}{\linewidth}{@{} l l Y Y @{}}
  \toprule
  \textbf{Op} & \textbf{Command} & \textbf{Payload} & \textbf{Reply} \\
  \midrule
    \texttt{0x01} & \texttt{WRITE}  & \texttt{addr} (input memory), \texttt{count}, \texttt{count} data bytes & none \\
  \texttt{0x02} & \texttt{READ}   & \texttt{addr} (result memory), \texttt{count} (elements), \texttt{size} (1 or 4 bytes per element)
                & \texttt{count}$\cdot$\texttt{size} bytes (\texttt{size} = 1: low byte, valid after \texttt{requant}) \\
  \texttt{0x03} & \texttt{MATMUL} & \texttt{aAddr}, \texttt{bAddr}, \texttt{cAddr}, \texttt{m}, \texttt{k}, \texttt{n},
                  \texttt{ldA}, \texttt{ldB}, \texttt{ldC}, flags, \texttt{shift}: 20 bytes, see Table~\ref{tab:matmul-frame}
                & none; poll \texttt{STATUS} \\
  \texttt{0x04} & \texttt{STATUS} & none & 9 bytes: flags (bit 0 busy, 1 error, 2 overrun), \texttt{cycles} (4), \texttt{busyPeCycles} (4) \\
  \bottomrule
\end{tabularx}

\begin{table}[h]
\small\renewcommand{\arraystretch}{1.2}
\begin{tabularx}{\linewidth}{@{} l l Y @{}}
  \toprule
  \textbf{Byte} & \textbf{Field} & \textbf{Meaning} \\
  \midrule
  0      & opcode & \texttt{0x03} \\
  1--2   & \texttt{aAddr} & input-memory address of $A[0][0]$ \\
  3--4   & \texttt{bAddr} & input-memory address of $B[0][0]$ \\
  5--6   & \texttt{cAddr} & result-memory address of $C[0][0]$ \\
  7--8   & \texttt{m}     & rows of A and C; $m \le \texttt{maxM}$ \\
  9--10  & \texttt{k}     & columns of A, rows of B; padded to a multiple of \texttt{banks} \\
  11--12 & \texttt{n}     & columns of B and C; padded to a multiple of \texttt{banks} \\
  13--14 & \texttt{ldA}   & distance between two rows of A in the input memory ($= k$ if dense) \\
  15--16 & \texttt{ldB}   & distance between two rows of B in the input memory ($= n$ if dense) \\
  17--18 & \texttt{ldC}   & distance between two rows of C in the result memory ($= n$ if dense) \\
  19     & flags          & bit 0 \texttt{requant}, bit 1 \texttt{relu}, bits 2--7 = 0 \\
  20     & \texttt{shift} & requantization shift ($0$ to $\texttt{accWidth}-1$); ignored if \texttt{requant} = 0 \\
  \bottomrule
\end{tabularx}
\caption{\texttt{MATMUL} frame: 1 opcode byte and 20 payload bytes, 2-byte fields little-endian.}
\label{tab:matmul-frame}
\end{table}

An unknown opcode is ignored and sets \texttt{error}. If more than $2^{20}$ cycles ($\approx$10\,ms) pass between two bytes of a frame, the partial frame is dropped, \texttt{error} is set and the decoder waits for an opcode again.

\paragraph{Status flags.}
\texttt{error} and \texttt{overrun} are sticky between reads: each is set by its event and cleared
when a \texttt{STATUS} reply is captured, so a reply reports what happened since the previous
\texttt{STATUS}. A new event in the same cycle as the capture wins over the clear, so no event is
lost. \texttt{busy} is a live level, never cleared. Reset clears both flags.
When \texttt{error} or \texttt{overrun} is set, an earlier frame may have been misread: the host
resends it. Recommended host habit: send \texttt{STATUS} after each \texttt{WRITE} and before each
\texttt{READ}, so that a flag always points to the last transfer.

Functions \texttt{quantize} and \texttt{dequantize} live in the Python program (and have a copy in Scala with the same rounding rule). Frédéric also chooses scales and \texttt{shift}.

% ------------------------------------------------------------
\subsection*{\texttt{UnifiedBuffer} (Giovanni)}
{\small\renewcommand{\arraystretch}{1.25}
\noindent\begin{tabularx}{\linewidth}{@{} l l l l l Y @{}}
  \porthead
  \texttt{clock} & in & Clock & -- & \texttt{Tpu} & 100\,MHz clock shared by all blocks. \\
  \texttt{reset} & in & Bool & W & \texttt{Tpu} & Synchronous, active-high. Pending reads are dropped. \textbf{Memory contents are not cleared}: matrices must be written again. \\
  \midrule
  \texttt{hostWr}            & in  & \texttt{BufWritePort}    & RV & \texttt{CommandDecoder} & Element writes from the host into the input memory. \texttt{ready} is constantly 1. \\
  \texttt{hostRdReq}         & in  & \texttt{BufReadReq}      & RV & \texttt{CommandDecoder} & Result read requests from the host. \\
  \texttt{hostRdResp}        & out & \texttt{accT}            & RV & \texttt{CommandDecoder} & Result values, one cycle after the request at the earliest. \\
  \texttt{mvRdReq}$\dagger$  & in  & \texttt{BufRowReq}       & RV & \texttt{DataMover}      & Row address: one value per bank is read in the same cycle. \\
  \texttt{mvRdResp}$\dagger$ & out & Vec(banks, \texttt{inT}) & RV & \texttt{DataMover}      & The row of weights or activations. \\
  \texttt{resWr}$\dagger$    & in  & \texttt{BufResultWrite}  & RV & \texttt{Accumulators}   & Final results of a tile, written back. \\
  \bottomrule
\end{tabularx}}

\paragraph{Internal organisation of \texttt{UnifiedBuffer}.}
Two independent memories, so no port is ever shared and no arbitration is needed:
\begin{itemize}
  \item \emph{Input memory} (\texttt{inT}): A and B. Written element by element by the host (\texttt{hostWr}, always \texttt{ready}), read one whole row per cycle by the data mover.
  \item \emph{Result memory} (\texttt{accT} words): C. Written one whole row per cycle by the accumulators (\texttt{resWr}), read element by element by the host (\texttt{hostRdReq}/\texttt{hostRdResp}).
  
\end{itemize}

The buffer does not check conflicts. The host must not \texttt{WRITE} addresses read by a running \texttt{MATMUL}, nor \texttt{READ} its results before \texttt{busy} drops. Safe sequence: \texttt{WRITE}s, \texttt{MATMUL}, poll \texttt{STATUS} until not busy, \texttt{READ}s. \\
\emph{Post-MVP:} to chain layers on-chip, requantized int8 results could be written into the input memory,
so that the output of one \texttt{MATMUL} becomes the A of the next, laid out by \texttt{ldC}.

% ------------------------------------------------------------
\subsection*{\texttt{Sequencer} with tiling controller (Giovanni)}
{\small\renewcommand{\arraystretch}{1.25}
\noindent\begin{tabularx}{\linewidth}{@{} l l l l l Y @{}}
  \porthead
  \texttt{clock} & in & Clock & -- & \texttt{Tpu} & 100\,MHz clock shared by all blocks. \\
  \texttt{reset} & in & Bool & W & \texttt{Tpu} & Synchronous, active-high. A running instruction is aborted; \texttt{busy} = 0, \texttt{instr.ready} = 1; counters set to 0. \\
  \midrule
  \texttt{instr}             & in  & \texttt{Instruction} & RV & \texttt{CommandDecoder} & Compute instruction. \texttt{ready} stays low until the previous one has finished. \\
  \texttt{status}            & out & \texttt{SeqStatus}   & W  & \texttt{CommandDecoder} & Busy flag and performance counters. \\
  \texttt{tileCmd}$\dagger$  & out & \texttt{TileCmd}     & RV & \texttt{DataMover}      & Which weight tile to load and which activation rows to stream. \\
  \texttt{accCtrl}$\dagger$  & out & \texttt{AccCtrl}     & RV & \texttt{Accumulators}   & First K-tile (overwrite) or not (add), last K-tile (write back), output address. \\
  \texttt{tileDone}$\dagger$ & in  & Bool                 & V  & \texttt{Accumulators}   & Pulse: the current tile's results are complete. \\
  \bottomrule
\end{tabularx}}

\paragraph{Internal operations.}
Weight loading, draining and activation are not instructions: for every tile the sequencer performs them itself, as internal steps of the \texttt{MATMUL} it is running.

\paragraph{Counters.}
\begin{itemize}
    \item \texttt{cycles} is a 32 bit register that becomes 0 when an instruction is accepted and counts every cycle until busy is 1.
    \item \texttt{busyPeCycles} doesn't get "measured", you sum $m*rows*cols$ at each tile. With $ros*cols$ power of 2 it's just a shift and in the end its value is $m*k*n$.
\end{itemize}



% ------------------------------------------------------------
\subsection*{\texttt{DataMover} (Giovanni)}
{\small\renewcommand{\arraystretch}{1.25}
\noindent\begin{tabularx}{\linewidth}{@{} l l l l l Y @{}}
  \porthead
  \texttt{clock} & in & Clock & -- & \texttt{Tpu} & 100\,MHz clock shared by all blocks. \\
  \texttt{reset} & in & Bool & W & \texttt{Tpu} & Synchronous, active-high. A tile being streamed is abandoned; \texttt{wOut.valid} = \texttt{aOut.valid} = 0. \\
  \midrule
  \texttt{cmd}$\dagger$    & in  & \texttt{TileCmd}         & RV & \texttt{Sequencer}     & Tile job: weight and activation addresses, number of rows M. \\
  \texttt{rdReq}$\dagger$  & out & \texttt{BufRowReq}       & RV & \texttt{UnifiedBuffer} & Row read requests. \\
  \texttt{rdResp}$\dagger$ & in  & Vec(banks, \texttt{inT}) & RV & \texttt{UnifiedBuffer} & Rows read from the buffer. \\
  \texttt{wOut} & out & Vec(cols, \texttt{inT}) & V & \texttt{SystolicArray}
                & One row of the B tile per valid cycle, \textbf{last row first}; \texttt{rows} valid cycles load a tile. \\
  \texttt{aOut} & out & Vec(rows, \texttt{inT}) & V & \texttt{SystolicArray}
                & One row of A per valid cycle, not skewed; cycles with valid = 0 allowed between rows.
                  First row at least $\texttt{weightToActGap}+1$ cycles after the last weight row. \\
  \bottomrule
\end{tabularx}}

% ------------------------------------------------------------
\subsection*{\texttt{SystolicArray} with skew and de-skew (Andreea)}
{\small\renewcommand{\arraystretch}{1.25}
\noindent\begin{tabularx}{\linewidth}{@{} l l l l l Y @{}}
  \porthead
  \texttt{clock} & in & Clock & -- & \texttt{Tpu} & 100\,MHz clock shared by all blocks. \\
  \texttt{reset} & in & Bool & W & \texttt{Tpu} & Synchronous, active-high. All \texttt{valid} bits inside the array are cleared, so no output appears until new inputs arrive. Weights and partial-sum registers keep their values (they are overwritten before use). \\
  \midrule
  \texttt{wIn} & in  & Vec(cols, \texttt{inT})  & V & \texttt{DataMover}
               & One row of the B tile per valid cycle; each valid cycle pushes the loaded weights down by one
                 PE row. Rows arrive last row first, so after \texttt{rows} valid cycles PE row $r$ holds tile row $r$.
                 Weights do not move while valid = 0. \\
  \texttt{aIn} & in  & Vec(rows, \texttt{inT})  & V & \texttt{DataMover}
               & One row of A per valid cycle: element $r$ is $A[m][kt\cdot\texttt{rows}+r]$ and goes to PE row $r$.
                 Not skewed (the array skews internally); cycles with valid = 0 are ignored. \\
  \texttt{out} & out & Vec(cols, \texttt{accT}) & V & \texttt{Accumulators}
               & One row of partial sums per valid cycle, de-skewed, exactly \texttt{arrayLatency} cycles after its
                 \texttt{aIn} row, same order and same gaps. Sums over the current K-tile only. Must always be accepted. \\
  \bottomrule
\end{tabularx}}

\paragraph{Array timing contract.}
The three ports of \texttt{SystolicArrayIO} (file \texttt{tpu/ifaces/ArrayInterface.scala}) use protocol V:
a cycle with valid = 1 carries exactly one row, a cycle with valid = 0 carries nothing, and the array never
makes its neighbours wait. The data mover and the array follow these rules:
\begin{enumerate}
  \item \emph{Loading weights.} The data mover sends the \texttt{rows} rows of the B tile on \texttt{wIn},
        last row first. Each valid cycle pushes the loaded weights down by one PE row, so that afterwards
        PE row $r$ holds row $r$ of the tile. While \texttt{wIn.valid} = 0 the weights stay in place.
  \item \emph{From weights to activations.} If the last weight row is valid at cycle $t$, the first
        activation row may be valid at cycle $t + 1 + \texttt{weightToActGap}$ at the earliest
        ($\texttt{weightToActGap} = 0$ for now, to confirm with Andreea).
  \item \emph{Streaming activations.} The data mover sends the rows of A for the current K-tile on
        \texttt{aIn}, one per valid cycle, not skewed. Gaps (valid = 0) between rows are allowed:
        the valid bit travels through the array with the data.
  \item \emph{Results.} Row $m$ of partial sums leaves on \texttt{out} exactly \texttt{arrayLatency} cycles
        after row $m$ entered on \texttt{aIn}, in the same order and with the same gaps. It contains sums over
        the current K-tile only; the accumulators add them across K-tiles.
  \item \emph{No overlap (MVP).} \texttt{wIn} and \texttt{aIn} are never valid in the same cycle. The next
        tile's weights start only after the last result row of the current tile has left the array, i.e.\
        \texttt{arrayLatency} cycles after its last \texttt{aIn} row. Overlap needs weight double buffering
        (post-MVP).
\end{enumerate}

% ------------------------------------------------------------
\subsection*{\texttt{Accumulators} (Giovanni)}
{\small\renewcommand{\arraystretch}{1.25}
\noindent\begin{tabularx}{\linewidth}{@{} l l l l l Y @{}}
  \porthead
  \texttt{clock} & in & Clock & -- & \texttt{Tpu} & 100\,MHz clock shared by all blocks. \\
  \texttt{reset} & in & Bool & W & \texttt{Tpu} & Synchronous, active-high. Control state reset, \texttt{done} = 0. Stored partial sums are not cleared (the first K-tile overwrites them). \\
  \midrule
  \texttt{in}              & in  & Vec(cols, \texttt{accT}) & V  & \texttt{SystolicArray} & Partial-result rows. Every valid beat must be accepted (the array cannot wait). \\
  \texttt{ctrl}$\dagger$   & in  & \texttt{AccCtrl}         & RV & \texttt{Sequencer}     & Overwrite or add, write back or keep, output address. \\
  \texttt{resWr}$\dagger$  & out & \texttt{BufResultWrite}  & RV & \texttt{UnifiedBuffer} & Final results after the last K-tile. \\
  \texttt{done}$\dagger$   & out & Bool                     & V  & \texttt{Sequencer}     & Pulse: all rows of the current tile received. \\
  \bottomrule
\end{tabularx}}
\paragraph{Post-processing.}
Applied to every element on the last K-tile only, in this order ($\gg$ is an arithmetic shift right):
\begin{enumerate}
  \item \texttt{requant} = 0: $y = \mathit{acc}$ (raw \texttt{accT}); \texttt{shift} is ignored.
  \item \texttt{requant} = 1: $y = \mathrm{sat}\big((\mathit{acc} + 2^{\mathit{shift}-1}) \gg \mathit{shift}\big)$, with rounding term 0 if \texttt{shift} = 0; the addition is done with one extra bit; \texttt{sat} clamps to $[-127, 127]$.
  \item \texttt{relu} = 1: $y = \max(y, 0)$.
\end{enumerate}
The result is always stored as an \texttt{accT}-wide word (with \texttt{requant} it lies in $[-127,127]$, or $[0,127]$ with ReLU).
\texttt{READ} chooses how many bytes per element travel to the host.

% ------------------------------------------------------------
\subsection*{Shared bundles}
Defined once in the \texttt{ifaces} package.

\begin{landscape}
{\small\renewcommand{\arraystretch}{1.3}
\begin{xltabular}{\linewidth}{@{} >{\raggedright\arraybackslash}p{3.3cm} >{\raggedright\arraybackslash}p{6.2cm} Y >{\raggedright\arraybackslash}p{3.4cm} @{}}
  \caption{Shared bundles. Bundles marked $\dagger$ are internal to Giovanni's part; their exact fields are his choice.}\label{tab:bundles}\\
  \toprule
  \textbf{Bundle} & \textbf{What it does} & \textbf{Fields} & \textbf{Used by} \\
  \midrule
  \endfirsthead
  \multicolumn{4}{@{}l}{\emph{(continued)}}\\
  \toprule
  \textbf{Bundle} & \textbf{What it does} & \textbf{Fields} & \textbf{Used by} \\
  \midrule
  \endhead
  \bottomrule
  \endlastfoot

  \multicolumn{4}{@{}l}{\textit{File \texttt{tpu/ifaces/HostInterfaces.scala}: host side}} \\[2pt]
  \texttt{BufWritePort}
    & Carries one element of A or B that the host sends with \texttt{WRITE}, and where to store it in the input memory.
    & \texttt{addr}: UInt(addrW); \texttt{data}: \texttt{inT}
    & Decoder $\rightarrow$ input memory \\
  \texttt{BufReadReq}
    & Asks the result memory for one element of C, for a \texttt{READ}. The value comes back on a separate \texttt{Decoupled(accT)}.
    & \texttt{addr}: UInt(outAddrW)
    & Decoder $\rightarrow$ result memory \\
  \texttt{Instruction}
    & Tells the sequencer to compute $C = A \cdot B$: where A, B and C are, their sizes and row spacing, and whether to requantize. Accepting it starts the computation.
    & \texttt{op}: \texttt{Opcode} (only \texttt{Matmul}); \texttt{aAddr}, \texttt{bAddr}: UInt(addrW); \texttt{cAddr}: UInt(outAddrW);
      \texttt{m}, \texttt{k}, \texttt{n}, \texttt{ldA}, \texttt{ldB}, \texttt{ldC}: UInt(dimW); \texttt{requant}, \texttt{relu}: Bool;
      \texttt{shift}: UInt(shiftW)
    & Decoder $\rightarrow$ sequencer \\
  \texttt{SeqStatus}
    & Tells the decoder whether a computation is running and how long the last one took, so it can answer \texttt{STATUS}. Plain wires, always valid.
    & \texttt{busy}: Bool; \texttt{cycles}, \texttt{busyPeCycles}: UInt(32)
    & Sequencer $\rightarrow$ decoder \\
  \texttt{DebugInfo}
    & Drives the board LEDs, so a human can see at a glance whether the design is busy or has lost data. Flags stay on until reset.
    & \texttt{busy}, \texttt{overrun}, \texttt{error}: Bool
    & \texttt{Tpu} $\rightarrow$ \texttt{BoardTop} \\
  \midrule

  \multicolumn{4}{@{}l}{\textit{File \texttt{tpu/ifaces/ArrayInterface.scala}: systolic array}} \\[2pt]
  \texttt{SystolicArrayIO}
    & The three ports of the array: weights in at the top, activations in on the left, partial sums out at the bottom. Each valid cycle carries one row. Rules in the array timing contract.
    & \texttt{wIn}: Valid(Vec(cols, \texttt{inT})), in; \texttt{aIn}: Valid(Vec(rows, \texttt{inT})), in;
      \texttt{out}: Valid(Vec(cols, \texttt{accT})), out (directions seen from the array)
    & Data mover $\rightarrow$ array $\rightarrow$ accumulators \\
  \midrule

  \multicolumn{4}{@{}l}{\textit{Inside Giovanni's part}} \\[2pt]
  \texttt{TileCmd}$\dagger$
    & Tells the data mover which tile to feed: where its weights and activations start, their row spacing, and how many rows of A to stream.
    & \texttt{aAddr}, \texttt{bAddr}: UInt(addrW) (first element of the tile's activations / weights); \texttt{ldA}, \texttt{ldB}, \texttt{m}: UInt(dimW)\,*
    & Sequencer $\rightarrow$ data mover \\
  \texttt{BufRowReq}$\dagger$
    & Asks the input memory for one full row (one value per bank) in a single cycle.
    & \texttt{rowAddr}: UInt(rowAddrW)
    & Data mover $\rightarrow$ input memory \\
  \texttt{BufResultWrite}$\dagger$
    & Writes one finished row of C (one value per PE column) into the result memory in a single cycle.
    & \texttt{rowAddr}: UInt($\texttt{outAddrW} - \log_2\texttt{cols}$); \texttt{data}: Vec(cols, \texttt{accT})
    & Accumulators $\rightarrow$ result memory \\
  \texttt{AccCtrl}$\dagger$
    & Tells the accumulators what to do with the next tile's results: start a new sum or add to it, and whether this is the last K-tile, in which case they post-process and write C.
    & \texttt{first}, \texttt{last}: Bool; \texttt{cAddr}: UInt(outAddrW); \texttt{ldC}, \texttt{m}: UInt(dimW); \texttt{requant}, \texttt{relu}: Bool;
      \texttt{shift}: UInt(shiftW)
    & Sequencer $\rightarrow$ accumulators \\
\end{xltabular}}

\noindent * The data mover computes the address of weight row $r$ as $\texttt{bAddr} + r\cdot\texttt{ldB}$ and of activation
row $i$ as $\texttt{aAddr} + i\cdot\texttt{ldA}$. The sequencer computes \texttt{aAddr} and \texttt{bAddr} for each tile from the
instruction's addresses, and copies \texttt{ldA}, \texttt{ldB} (and \texttt{ldC} into \texttt{AccCtrl}) from the instruction.
\end{landscape}

\vspace{0.7em}

* The data mover computes the address or row $r$ of weights as $bAddr + r*ldB$ and the one of the activation $i$ as $aAddr + i*ldA$; the sequencer computes $aAddr$, $bAddr$ for each tile from addresses of the instruction.

\end{document}
